\section{Related Work}
\label{sec:related-work}

Classical UAV path planning methods include graph-search algorithms, sampling-based planners, and potential-field approaches. Graph-search methods such as Dijkstra and A* are widely used because they are interpretable and provide strong guarantees when the environment is represented as a known graph with non-negative edge weights. Sampling-based planners such as RRT and PRM can explore larger spaces, but their randomized nature means they do not necessarily return optimal paths and can struggle in cluttered environments \cite{ghambari2024survey,venu2025comprehensive}. This project uses Dijkstra rather than a sampling-based method as the classical baseline because it is deterministic, optimal on a static graph, and has well-characterized computational complexity, properties that make the trade-off being measured easier to state precisely.

Recent surveys of UAV path planning emphasize that classical planners remain important reference points because of their completeness, optimality, and predictable computational behavior \cite{mannan2023classical,ghambari2024survey,johnson2026rcsr}. Their main limitation in dynamic environments is that they generally assume a known graph. When obstacles or edge costs change, the planner must update or recompute a route.

Reinforcement learning has been explored as an alternative for UAV routing and obstacle avoidance. RL methods learn policies through interaction rather than explicit shortest-path search. Deep RL methods such as DQN are particularly relevant when the state space becomes too large for tabular Q-learning \cite{guo2023uav,tu2023uav,zhao2022path}. Prior work reports that RL can improve autonomy and adaptability in complex settings, but it often trades away the guarantees of classical methods.

Dynamic and online UAV routing is especially relevant to this project. Rocha et al. propose dynamic Q-planning for online UAV path planning in unknown and complex environments, comparing reinforcement-learning-based planning with classical alternatives \cite{rocha2024dynamic}. This supports the motivation for testing learning-based adaptation against a deterministic replanning baseline.

Much of the RL-versus-classical UAV literature compares methods across different metrics, environments, and hardware, making head-to-head comparison across papers difficult; few evaluate both approaches on the exact same graph, obstacle layout, and start-goal pairs with a paired statistical test on the difference. This project fits into the literature by closing that specific gap: a controlled same-scenario comparison rather than a broad benchmark. Both methods are evaluated on the same grid, same dynamic cells, and same 50 start-goal pairs. This paired design allows statistical testing of path-cost and compute-time differences instead of relying on anecdotal comparisons.
